\section{Measurement Protocol}
\label{sec:protocol}

\paragraph{How large is the noise?}
We ran the baseline four times on ImageNet-R with one seed. The base task
was identical in each run; only the SGD randomness of later tasks changed.
The resulting spread is $\bar A=87.35\pm0.14$ and $A_B=82.59\pm0.23$; on
ImageNet-A the spread of $\bar A$ is $0.38$. Changing the seed also changes
the class order and adds more variation. A single run of a candidate
against a single run of the baseline therefore cannot resolve effects below
roughly half a point.

\paragraph{Reproduction.}
Using the authors' code and data partitions, our runs reproduce all twelve
numbers of the baseline's main table ($\bar A$ and $A_B$ on six benchmarks) to
within $0.01$. Our experiments then use partitions we can redistribute:
\begin{itemize}
  \item the standard CIFAR-100 split;
  \item a public fixed ImageNet-R split with 23\,918 training and 6\,082 test images;
  \item an 80/20 per-class split of ImageNet-A, with seed 67;
  \item the official CUB-200 split.
\end{itemize}
Absolute numbers therefore differ slightly from the published ones.

The protocol has four rules, described next.

\paragraph{Rule 1: pair everything.}
Every dataset is run with the same three seeds, 1993, 2015 and 2026.
Effects are reported as paired differences within a run, with 95\%
$t$-intervals across runs.

\paragraph{Rule 2: evaluate all heads on the same network.}
A component that only changes inference does not need its own training run.
After each task, the linear head and every candidate head are evaluated on
the same test features. Each head draws its randomness from a private
generator, and the global random state is saved and restored around the
evaluation. Two checks are run every time:
\begin{itemize}
  \item the linear head reproduces the network's logged accuracy curve;
  \item that curve matches, digit for digit, the curve of a plain baseline
    run with the same seed and replicate.
\end{itemize}
Both held in all runs. Training noise then cancels in head-to-head
differences. On ImageNet-R, the standard deviation of a paired difference is
$0.03$--$0.05$ points instead of $0.2$, so three runs resolve $0.2$ points.

Components that change training, such as a quantum steering loss, cannot
use this rule. They are compared with baseline runs of the same seed and
replicate, at the larger noise level.

\paragraph{Rule 3: decide before looking.}
Each round of experiments had a written specification before it ran. It
fixed the components, the hyper-parameters, the decision rule and the
minimum effect. Free choices, such as fusion weights, were made once, on a
tuning group (ImageNet-R, seed 1993, three replicates), and then frozen.
CIFAR-100 and CUB-200 were added only after all choices had been made, and
serve as pure tests of transfer. Table~\ref{tab:tests} lists every test,
including those that failed.

\paragraph{Rule 4: a classical twin for every quantum component.}
A quantum component is credited only if it beats a classical method that
receives the same inputs and the same budget:
\begin{itemize}
  \item the PQK is compared with an RBF kernel or random features on the
    same projection;
  \item a circuit with entangling gates is compared with the same circuit
    without them;
  \item the PGM is compared with ridge regression on the same class
    statistics.
\end{itemize}
Beating the baseline alone is not enough. The baseline does not use the
extra information that the component and its twin share.
